\subsection{Distinguished $S$ symmetry over the full character field}
\label{sec:rank8-full-character-S}
For a commutative pseudo-unitary fusion ring, let $\chi_E$ run through its
irreducible characters, $f_E$ be the formal codegree, and
$\beta_E=D/f_E$. The distinguished central simple $A_E$ in the induction
of the tensor unit has dimension $\beta_E$ and multiplicity one. If
$v_E=[F(A_E)]$, the induction--character formula and symmetry of $S$ give
\begin{equation}\label{eq:rank8-full-character-S}
 \beta_E\chi_E(v_F)=\beta_F\chi_F(v_E).
\end{equation}
Here the formula is symmetric, without complex conjugation. It follows
from Deshpande, \emph{On centers of bimodule categories and
induction-restriction functors}, arXiv:1611.04486, Lemma~4.3 and
Corollary~4.5, specialized to the identity bimodule and one-dimensional
characters. The equations involving the Frobenius--Perron character
hold automatically because its distinguished simple is the tensor unit.

We extend the atomic induction relaxation by taking one simplex for
each non-Frobenius--Perron character, over all allowed integral
restriction columns of dimension $\beta_E$, and imposing
\eqref{eq:rank8-full-character-S} in every rational coordinate of the
full character splitting field $L$. We first apply the necessary
algebraic-integrality test to
\[
 \frac{\beta_E\chi_E(v)}{\operatorname{FPdim}(v)}.
\]
Integral closure is tested by the characteristic polynomial of the
rational multiplication operator on $L$. These symmetry equations
involve only distinguished columns: high columns and rotation columns
have zero coefficients in the added rows. Thus the twist-moment
computations remain in the original Frobenius--Perron field and no
extension of their relative-trace field is required. No $STS$ equation
is used in this variant. Independent simplex mixtures give a necessary
relaxation, weaker than the complete joint-frame hull; feasibility has
no realization implication.

For $R_{8,3,3413}$ the full field is
$\mathbb Q(\sqrt5,\theta)$, where $\theta^3-3\theta+1=0$.
Using $u=\sqrt5+\theta$ gives
\[
 u^6-21u^4+2u^3+84u^2+24u-19=0.
\]
The exact height-four cone has $253$ equations and $4485$ columns:
$780$ distinguished choices, $367$ high restriction vectors paired with
$10$ joint trace types, and $35$ rotation columns. The sole permitted
second-indicator assignment is $(1,1,1,1,1,1,1)$. The saved integral
functional is nonnegative on every column and strictly negative on
the right-hand side. Every inequality is also checked by direct column
dot products independently of the factored oracle.

The general lift has five spherical-character factors, of degrees
$1,1,1,2,3$. Four factors force a zero simple-object dimension; the
quadratic factor is the Frobenius--Perron Galois orbit, excluded by the
preceding cone. The complete oriented signed enumeration has one
branch, gauge equivalent to the original ring. Therefore
$R_{8,3,3413}$ admits no complex fusion-category categorification,
without a pivotality or unitarity assumption. The independent replay
is \texttt{quadratic/review\_atomic\_spectrum.py}; its source-bound
receipt also includes a known $\mathrm{Fib}\boxtimes\mathrm{Vec}_{C_2\times C_2}$
exact feasible control and complete oracle comparisons.
